\chapter{Verifying Candidate Solutions}\label{chap:4}

\begin{quote}
\emph{``Show me your flowchart and conceal your tables, and I shall
continue to be mystified. Show me your tables, and I won't usually need
your flowchart; it'll be obvious.''} --- Fred Brooks, \emph{The Mythical
Man-Month}
\end{quote}

In this chapter, we discuss a concept that is orthogonal to the basic
modeling process. A \emph{verifier} is a piece of code that takes an
instance and a candidate solution, checks whether the solution is
feasible, and computes its objective value.

Writing such a verifier still requires the same five ingredients from
optimization modeling: entities, parameters, variables, constraints, and
objective. In the verifier, however, the variables appear as concrete
values inside a candidate \lstinline!Solution!, not as
symbolic objects controlled by a solver. The difference is the medium.
Instead of expressing the problem as mathematical notation or as solver
constraints, we express it as ordinary program code. For many software
engineers, this feels less abstract and more natural.

A verifier can be developed independently of the solver model, but the
two must agree on the problem definition and on the representation of
instances and solutions. Once available, the verifier can serve several
roles. It can help you test hand-written examples, evaluate simple
heuristics, and get a feel for the problem before the solver model is
mature. During solver development, it acts as a guardrail:
solver-produced solutions can be checked against it, random instances
can be smoke-tested, and performance-driven refactors can be validated
against the same problem definition.

In that sense, the verifier can be viewed as an executable contract for
the problem definition. It does not say how to find a good solution. It
says how to recognize a valid one, and how to score it once found.

\section{What a Verifier Is}\label{sec:4.1}

A verifier checks candidate solutions. Given an instance and a proposed
solution, it answers two questions:

\begin{enumerate}
\item
  Is the solution feasible?
\item
  What is its objective value?
\end{enumerate}

It does not search for solutions. It does not call a solver. It is not
an algebraic modeling language and not a solver model. A solver model
describes a search space to an optimization engine; a verifier checks
one fully specified candidate after the decisions have already been
made. Solvers produce solutions; verifiers consume them.

This difference makes a verifier much easier to write for many rules. In
a solver model, a rule must be expressed as a constraint over decision
variables, in the vocabulary supported by the solver. In a verifier, the
candidate assignment is already known, so the rule can be checked with
ordinary program code: branches, lookups, loops, early returns, and
descriptive error messages.

The \textbf{simple form} keeps all checks in one evaluator function. It
returns the objective value if the candidate is feasible and raises an
exception naming the first violated rule otherwise. This is short,
readable, and sufficient for many small hard-constraint models.

The \textbf{richer form} decomposes the checks into separate rule
modules and replaces yes/no answers with quantitative metrics:
\emph{three minutes of overlap}, \emph{two uncovered shifts}, \emph{one
missing skill}. This becomes useful once rules are soft, violations need
to be reported in bulk, or several solution methods need to be compared
against the same standard.

\section{\texorpdfstring{The Simple Form: \texttt{Instance}, \texttt{Solution}, \texttt{evaluate}}{The Simple Form: Instance, Solution, evaluate}}\label{sec:4.2}

A simple verifier consists of three pieces. The first two are plain data
containers and can be reused elsewhere in the codebase:

\begin{enumerate}
\item
  \textbf{\lstinline!Instance!}, a frozen dataclass
  holding the input data.
\item
  \textbf{\lstinline!Solution!}, a frozen dataclass
  holding one completed candidate solution.
\item
  \textbf{\lstinline!evaluate(instance, solution)!}, a
  pure function that checks feasibility and, if the solution is
  feasible, returns its objective value. If a constraint family fails,
  it raises an exception that identifies the failing rule.
\end{enumerate}

This is the minimal pattern: two dataclasses and one function. We illustrate it with a rostering verifier.

\subsection{Type Aliases for the Index Sets}\label{sec:4.2.1}

Python is dynamically typed, but type annotations are useful as
documentation and for static analysis and editor support. Here, we use
simple aliases for the index sets from the paper model:

\begin{lstlisting}[language=Python]
type Worker = str  # Identifier for a worker who can take shifts
type Shift  = str  # Identifier for a shift that needs to be staffed
type Day    = int  # Integer day index, e.g., 0 for Monday, 1 for Tuesday
type Skill  = str  # Skill that workers may hold and shifts may require
\end{lstlisting}

Workers, shifts, days, and skills are still strings or integers at
runtime. The aliases do not add type safety by themselves, but they
improve the shape of the code:
\lstinline!dict[tuple[Worker, Shift], int]! says more than
\lstinline!dict[tuple[str, str], int]!.

\subsection{\texorpdfstring{The \texttt{Instance} and \texttt{Solution} dataclasses}{The Instance and Solution dataclasses}}\label{sec:4.2.2}

\begin{lstlisting}[language=Python]
from dataclasses import dataclass

@dataclass(frozen=True)
class Instance:
    workers: frozenset[Worker]
    shifts:  frozenset[Shift]
    days:    frozenset[Day]
    skills:  frozenset[Skill]

    cost:           dict[tuple[Worker, Shift], int]
    available:      dict[tuple[Worker, Shift], bool]
    required_skill: dict[Shift,  Skill]
    holds_skills:   dict[Worker, frozenset[Skill]]
    headcount:      dict[Shift,  int]
    shifts_on_day:  dict[Day,    frozenset[Shift]]

@dataclass(frozen=True)
class Solution:
    assigned: frozenset[tuple[Worker, Shift]]
\end{lstlisting}

The \lstinline!Instance! carries the sets and parameters
from the math model; the \lstinline!Solution! carries the
values of the decision variables. The solution uses the \emph{sparse}
representation (we record only the pairs where \(x_{w,s} = 1\), not
the full zero-one matrix) which scales with what was assigned rather
than with the size of the index space.

\begin{platypusinfo}
For teaching purposes we use
\lstinline!dataclass!, but at the system boundary,
wherever an instance is built from a CSV, a database, or a JSON payload,
a properly validated data schema (typically
\href{https://docs.pydantic.dev/}{Pydantic}) earns its keep.
Inconsistent input data is extremely common in practice, and a
surprising share of ``the model is wrong'' debugging sessions turn out
to be data problems wearing a model's clothes. The topic deserves its
own chapter and we will not develop it here; what matters for the
verifier is only that the object handed to
\lstinline!evaluate! has the right attributes.
\end{platypusinfo}

\subsection{The Evaluator}\label{sec:4.2.3}

The evaluator is where each constraint family receives its explicit
check, and where the objective is computed only after all checks have
passed.

\begin{lstlisting}[language=Python]
def evaluate(inst: Instance, sol: Solution) -> int:
    """
    Return the objective value, or raise ValueError naming the failing rule.
    """

    # (cover) every shift gets exactly the headcount it requires
    for s in inst.shifts:
        got = sum((w, s) in sol.assigned for w in inst.workers)
        if got != inst.headcount[s]:
            raise ValueError(
                f"(cover) shift {s!r}: need {inst.headcount[s]}, got {got}"
            )

    # (avail) no assignment to an unavailable (w, s) pair
    for (w, s) in sol.assigned:
        if not inst.available[w, s]:
            raise ValueError(
                f"(avail) worker {w!r} not available for shift {s!r}"
            )

    # (skill) every assigned worker holds the shift's required skill
    for (w, s) in sol.assigned:
        required = inst.required_skill[s]
        if required not in inst.holds_skills[w]:
            raise ValueError(
                f"(skill) worker {w!r} lacks {required!r} required by shift {s!r}"
            )

    # (no-double) no worker is on two shifts on the same day
    for w in inst.workers:
        for t, shifts_today in inst.shifts_on_day.items():
            taken_today = sum((w, s) in sol.assigned for s in shifts_today)
            if taken_today > 1:
                raise ValueError(
                    f"(no-double) worker {w!r} on {taken_today} shifts on day {t}"
                )

    # objective: total cost over assigned pairs
    return sum(inst.cost[w, s] for (w, s) in sol.assigned)
\end{lstlisting}

The exception messages name witnesses, not just rules. For example,
\lstinline!(skill) worker 'Alice' lacks 'ICU' required by shift 'S5'!
identifies the violated rule together with the worker and shift
involved. A domain expert can read this message and confirm or dispute
it without understanding the solver.

\begin{platypustip}
Quantitative verifier output is also useful for AI-assisted
development. A coding agent can use violations, costs, and witness
messages as a concrete feedback signal when implementing a heuristic,
repairing a solution, or debugging a solver interface. The verifier does
not make the agent reliable by itself, but it gives the agent something
precise to test against.
\end{platypustip}

\section{A Richer Form: Quantitative Metrics and Modularized Rules}\label{sec:4.3}

On larger problems, two refactors of the simple form repeatedly pay off.
The exact shape of a richer verifier is project-specific; what follows
is the direction, not a blueprint.

\subsection{Quantitative Metrics Instead of Raise/Continue}\label{quantitative-metrics-instead-of-raisecontinue}

The simple \lstinline!evaluate! function answers a binary
question per rule: did the rule pass, or did it reject the solution?
That throws away information that becomes essential once a project has
\emph{soft} rules (rules that may be violated at a penalty rather
than forbidden outright) or any kind of repair-style search. Three
minutes of overlap and three hours of overlap are very different
solutions, but the simple form cannot distinguish them.

Replacing the raise/continue verdict with a
\lstinline!Metric! object preserves that information:

\begin{lstlisting}[language=Python]
@dataclass(frozen=True)
class Metric:
    rule: str
    violation: float            # 0.0 = rule satisfied; otherwise rule-native units
    cost: float                 # contribution to the objective
    messages: tuple[str, ...]   # one per witness, e.g. "shift 'S3': need 3, got 2"
\end{lstlisting}

The \lstinline!violation! field measures how far the
solution is from satisfying the rule, in whatever unit is natural for
that rule: uncovered shifts, minutes of overlap, or hours over the
weekly ceiling. The \lstinline!cost! field measures the
rule's contribution to the objective.

For hard rules, feasibility is decided by
\lstinline!violation!; a solution is feasible exactly when
total hard-rule violation is zero. For soft rules, the violation
magnitude often becomes part of the cost instead of making the solution
infeasible.

The \lstinline!messages! field carries the witness-naming
we already used in the simple form's exception strings (for example,
\lstinline!(cover) shift 'S3': need 3, got 2!) but now
per witness rather than only for the first failure. A debugging session
can therefore see every place a rule failed instead of only the point at
which \lstinline!evaluate! raised.

The gain is that solutions can now be compared by how badly they fail a
rule, which is what makes repair, local search, and soft-objective
optimization practical.

\begin{platypustip}
Given such feedback, coding agents can frequently implement
or fix algorithms autonomously as it provides a clear signal of whether
it is getting closer or not.
\end{platypustip}

\subsection{Splitting Rules into Modules}\label{splitting-rules-into-modules}

The other refactor is structural. The simple
\lstinline!evaluate! puts every rule in one function. With
four rules, that is fine; with twenty, it becomes hard to navigate, hard
to test in isolation, and hard to give per-rule precomputation a natural
place to live.

Splitting the rules across modules (typically as classes, grouped
however makes sense for the project, such as one file per rule or one
file per related family of rules) addresses all three issues. The
example below shows the pattern with a single rule per class:

\begin{lstlisting}[language=Python]
# verifier/rules/cover.py

class CoverRule:
    """Each shift must be staffed with exactly its required headcount."""

    def __init__(self, instance: Instance):
        self.instance = instance
        # Precompute whatever this rule needs per instance.

    def evaluate(self, solution: Solution) -> Metric:
        violation = 0.0
        messages: list[str] = []

        for s in self.instance.shifts:
            got = sum((w, s) in solution.assigned for w in self.instance.workers)
            need = self.instance.headcount[s]

            if got != need:
                violation += abs(got - need)
                messages.append(f"shift {s!r}: need {need}, got {got}")

        return Metric(
            rule="cover",
            violation=violation,
            cost=0.0,
            messages=tuple(messages),
        )
\end{lstlisting}

Each rule gets a constructor, which runs once per instance, and an
\lstinline!evaluate! method, which the top-level verifier
may call many times while checking candidates from a solver, heuristic,
repair procedure, or test suite. Anything that depends only on the
instance can be precomputed in the constructor. How the per-rule metrics
are aggregated into a final report is a project-specific decision and is
not worth prescribing here.
